\paragraph{Фаза~0: міст промпт--коміт.}
Перед розгортанням повноцінного рівня оператора встановлюється мінімальний колектор даних.
Хук \texttt{UserPromptSubmit} у Claude Code створює orphan"=коміт у bare"=репозиторії git для кожного промпту оператора.
Кожний запис містить: текст промпту, мітку часу, ідентифікатор сесії, репозиторій, гілку та режим дозволів.

Протягом 4--6~тижнів пасивного збору формується природномовний корпус, що забезпечує вхідні дані для трьох задач:
\begin{itemize}[nosep]
  \item розмежувальні мітки (disambiguation labels) --- класифікація типу наміру оператора;
  \item розподіли частот тем (topic"=frequency distributions) --- виявлення домінантних робочих контекстів;
  \item патерни перемикання між проєктами (cross"=project switching patterns) --- моделювання переходів оператора між репозиторіями та завданнями.
\end{itemize}

\noindent
Цей корпус є сировинним матеріалом для embedding"=pipeline рівня оператора: з нього витягуються структуровані абстракції, описані вище, без збереження сирого контенту промптів у довготривалій пам'яті.

Продуктивність: затримка захоплення становить 15\,мс на один промпт --- непомітна в межах 5"=секундного таймауту хуків Claude Code.
Фаза~0 розгорнута з 9~травня 2026~р.

\medskip
\noindent\textit{Військова аналогія.}
Фаза~0 аналогічна пасивному збору радіоелектронної розвідки (SIGINT): накопичення даних без втручання в оперативну діяльність, формування базового розуміння патернів поведінки оператора до розгортання активних систем аналізу та адаптації.
